%======================================================================
%  Main figures and tables
%
%  Figures 1, 2 and 4 are the assembled composites from the three companion
%  artifacts; the per-panel "component" cards those artifacts also display
%  are deliberately not reproduced (see CONVERSION_NOTES.md). Figure 3 is
%  reserved for the unwritten functional-analysis section.
%
%  Tables 1-2 are main-text tables. The origin artifact's third table, the
%  32-row parameter longtable, moved to Extended Data Table 1 on 2026-09-04 to
%  bring the main display count to the Nature Biotechnology limit of six.
%
%  \fitfig caps height as well as width so no figure runs off the page.
%======================================================================

% ---- Figure 1: pretraining ------------------------------------------
% Asset: figures/Fig_1_Main.pdf  (artifact/figures/pdf_files/Fig_1_Main.pdf,
% the supplied Illustrator composite, copied unchanged)
%
% ACTION FOR AUTHOR (2026-09-04): panel d of the ARTWORK still reads "~160k"
% where the text and this caption both give the measured 166,192. The caption
% used to document that disagreement in parentheses, which is not something a
% published caption can do, so the parenthetical was removed. The artwork now
% has to be corrected to match, or the number changed in both places.
\begin{figure}[h]
    \centering
    \fitfig{figures/Fig_1_Main}
\caption{
\textbf{\model: a genomic language model of plasmid sequence space.} \textbf{a}, Origin of transfer (\emph{oriT}): a relaxase nicks the plasmid at \emph{oriT} and the single strand is transferred through the conjugation pilus from donor to recipient cell. \textbf{b}, Origin of replication (\emph{oriV}): a Rep protein binds \emph{oriV} and initiates plasmid replication. \textbf{c}, GOLD ecosystem composition of the 693,638-sequence training corpus, shown as a treemap with tile area proportional to sequence count, alongside the reserved host-phylogeny area \cite{mukherjee2023twenty}. \textbf{d}, PTU-category composition of the 213,393 PTUs, 500,466 dereplicated sequences, and a mean 166,192 representatives per epoch; ribbons connect the same categories, and Major PTUs are split by dereplicated sequences per PTU. \textbf{e}, Per-PTU sampling rate relative to one sequence per PTU, measured after dereplication: sequence-uniform sampling reaches 715$\times$, whereas diversity-aware sampling is capped at 4.0$\times$. \textbf{f}, Hierarchical epoch resampling across PTUs, sub-PTUs, dereplication clusters, and sequence representatives. \textbf{g}, Masked language modelling with a 2,048 bp input window and a 72-layer ByteNet encoder, yielding per-nucleotide embeddings and nucleotide probabilities.
}
\label{fig:pretraining}
\end{figure}


% ---- Figure 2: origin classification ---------------------------------
% Asset: figures/Fig_2_Main.pdf -- a re-laid-out export supplied by the author
% (commit "Better sized fig 2"), replacing the origin artifact's own assembled
% card, figure4_composite.pdf. This is the one figure not copied unchanged
% from its source artifact; see CONVERSION_NOTES.md.
\begin{figure}[h]
    \centering
    \fitfig{figures/Fig_2_Main}
\caption{
\textbf{\model identifies origin-containing regions in plasmids.} \textbf{a}, Similarity-aware partitioning of the curated origins (112 \emph{oriT}, 224 \emph{oriV}) into five parts, and the cross-validation setup. \textbf{b}, Window-level detection AUPRC on the evaluated intergenic-window population, all methods (\emph{oriT} over \emph{oriV}). \textbf{c}, Window AUPRC on origins novel to training versus origins with training similarity $>$20 --- the operational generalization test (\emph{oriT} over \emph{oriV}). \textbf{d}, Region-level ranking of intergenic regions by Precision@1 (\emph{oriT} over \emph{oriV}). \textbf{e}, Calibrated window-probability tracks along two example regions. \textbf{f}, Per-region inference speed for a single-split, five-run window predictor. Stacked pairs in \textbf{b}--\textbf{d} read \emph{oriT} over \emph{oriV}; the model legend sits under \textbf{c}.
}
\label{fig:origin}
\end{figure}


% ---- Figure 3: functional analysis of discovered origins (RESERVED) ----
% PLACEHOLDER -- to be written/uploaded by author. This float exists so that
% the interpretation composite auto-numbers as Figure 4. Section retitled
% 2026-09-04, author-directed.
\begin{figure}[h]
    \centering
    \placeholderfig{6cm}{figures/functional analysis of discovered origins --- to be supplied by author}
\caption{
\todo{Functional analysis of \model-discovered origins --- figure to be supplied by author.}
}
\label{fig:genomewide}
\end{figure}


% ---- Figure 4: model interpretation ----------------------------------
% Asset: figures/fig5_catjac.pdf -- the artifact's "Final" assembled card.
%
% 2026-09-04: the caption's last two sentences were removed. The coordinate
% verification method belongs in Methods, where it already appears under
% "Locus curation and coordinate verification"; the internal checkpoint
% identifier v3_noflex_375k should not appear in a published caption. The
% released checkpoint is identified in Methods and in Code Availability.
% (Filename still says fig5 from an earlier numbering.) The "Fig. 4a" and
% "Fig. 4b" standalone component panels in that artifact are not reproduced.
\begin{figure}[h]
    \centering
    \fitfig{figures/fig5_catjac}
\caption{
\textbf{Model-derived dependency patterns at a characterized origin of transfer and origin of replication.} Categorical-Jacobian dependency maps computed with frozen \model over a 2,048 bp context and cropped to the locus (heatmap; each panel scaled to its own crop with the upper limit at the crop's 99th percentile, so intensities are not comparable between panels). Above each map: a logo track of per-position information content from the model's masked predictive distribution, measuring predictive specificity against a uniform background, and annotation lanes carrying externally sourced features as arrows (blue, inverted repeats; orange, direct repeats; green, promoter boxes). Red boxes are annotation-driven overlays marking the extent of an annotated feature, not model-detected boundaries. Axes are labelled in 1-based coordinates; intervals below are 0-based half-open. \textbf{a}, Origin of transfer of \emph{Klebsiella pneumoniae} plasmid pIGRK (AY543071.1, [904,1214)). IR1--IR3 and the two identical 12 bp direct repeats DR1, DR2 are the repeats proposed in the source study, which mapped the minimal \emph{oriT} interval experimentally and predicted the individual elements; $-$35 and $-$10 mark a P\emph{mobK} promoter predicted computationally in that study, not one verified experimentally, and DR2 lies within their spacer. IR4 ([1119,1152)) is not annotated in the source literature and was identified after inspecting this map. \textbf{b}, Minimal origin of replication of \emph{Escherichia coli} plasmid pSC101 (NC\_002056.1, [5100,5272)), a crop within the reported minimal origin: the DnaA box and IHF site lie outside the displayed interval. DR1--DR3 are the three 21 bp RepA iterons; IR1 and IR2 are the two inverted repeats, with the $-$35 and $-$10 boxes of the \emph{rep} promoter inside IR1 and IR2 respectively.
}
\label{fig:interpretation}
\end{figure}


\clearpage

% ---- Table 1 ---------------------------------------------------------
\begin{table}[h]
\caption{
\textbf{Region-level ranking, mean average precision (mAP)} averaged equally over the five similarity-aware splits. \emph{oriT} and \emph{oriV} evaluated separately; higher is better. mAP is per-plasmid average precision, macro-averaged within each split. Precision@1 is shown in \fref{fig:origin}d. Recall@1, NDCG and the per-split spread of every method are in \supptref{tab:supp-region-metrics}; P@1 and R@1 are not identical when a plasmid has more than one positive region. The reference-tool database ablation is described in \nameref{sec:methods}.
}
\label{tab:region-map}
\begin{center}
\small
\begin{tabular}{p{0.52\textwidth}cc}
\toprule
Method & \emph{oriT} mAP & \emph{oriV} mAP \\
\midrule
\multicolumn{3}{l}{\textit{Baselines}} \\
Longest intergenic (naive) & 0.352 & 0.663 \\
Reference tool (train) --- oriTfinder2 / OriV-Finder & 0.149 & 0.730 \\
Reference tool (train + val) & 0.157 & 0.732 \\
One-hot, CNN (from scratch) & 0.509 & 0.686 \\
One-hot, Residual CNN & 0.493 & 0.671 \\
\midrule
\multicolumn{3}{l}{\textit{Frozen DNA language-model embeddings}} \\
PlasmidGPT & 0.438 & 0.531 \\
Evo2-1B & 0.848 & 0.807 \\
Evo2-7B & 0.883 & 0.825 \\
Evo2-40B & 0.883 & 0.810 \\
\model (standalone, deployment) & 0.897 & 0.816 \\
Ensemble (\model + Evo2-7B) & 0.928 & 0.861 \\
\bottomrule
\end{tabular}
\end{center}
\end{table}


% ---- Table 2 ---------------------------------------------------------
\begin{table}[h]
\caption{
\textbf{Window-level detection, average precision (AUPRC)} on held-out test windows (mean $\pm$ s.d. over the five splits). The random-ranker AUPRC equals the positive-window prevalence, 0.048. Validation-set values are in \supptref{tab:supp-window-auprc}.
}
\label{tab:window-auprc}
\begin{center}
\begin{tabular}{lcc}
\toprule
Method & \emph{oriT} AUPRC & \emph{oriV} AUPRC \\
\midrule
One-hot, CNN & 0.277 $\pm$ 0.077 & 0.383 $\pm$ 0.076 \\
One-hot, Residual CNN & 0.279 $\pm$ 0.103 & 0.375 $\pm$ 0.030 \\
PlasmidGPT & 0.150 $\pm$ 0.064 & 0.201 $\pm$ 0.062 \\
Evo2-1B & 0.605 $\pm$ 0.148 & 0.570 $\pm$ 0.046 \\
Evo2-7B & 0.722 $\pm$ 0.100 & 0.616 $\pm$ 0.090 \\
Evo2-40B & 0.696 $\pm$ 0.119 & 0.546 $\pm$ 0.071 \\
\model (standalone) & 0.781 $\pm$ 0.150 & 0.671 $\pm$ 0.074 \\
Ensemble (\model + Evo2-7B) & 0.814 $\pm$ 0.104 & 0.692 $\pm$ 0.078 \\
\bottomrule
\end{tabular}
\end{center}
\end{table}
